% !TeX root = ../nc-quadric-classification.tex
% 2026-11-21（証明をBへ移管）
\section{一般 cubic helix の持上げと cubic locus の滑らかな開性}
\label{sec:cubic-smooth-atlas}

\paragraph{主張の範囲と依存関係.}
本節では、一般 geometric cubic 点について、De Deken--Lowen の helix 持上げから
四つの隣接 Hom 空間と四次 syzygy の直線を回復する。
可換二次曲面上の Künneth 計算は用いず、指定された pointed 同型と Hom の基底を保って
五つ組の有限パラメータ空間へ持ち上げる。
相対 Picard 根と定理~\ref{thm:geometric-cubic-relative-effectivity} による実際の族を入力とし、
幾何点がちょうど cubic 対象である有限型で滑らかな Artin 開部分を得る。
noncubic 境界の相対認識と $\mathfrak H_2$ 全体の代数性は別の問題である。
本節の相対回復と atlas 構成の新規性は未確認である。

\subsection{五つ組と四項 foundation}

$k$ は標数零の代数閉体とする。
四つの二次元 $k$ 空間 $U_0,U_1,U_2,U_3$ を固定し、
\[
 S\subset\mathscr G\subset
 \mathbb P(U_0\otimes_kU_1\otimes_kU_2\otimes_kU_3)
\]
を geometric 五つ組の開部分 $\mathscr G$ の任意の開部分とする。
ここで $\mathbb P$ は直線をパラメータ化する規約であり、
$S$ 上の tautological 直線を $W$ と書く。
非線型 locus に制限する場合も、以下の持上げ補題はそのまま適用できる。

$S$ 上で縮約により
\[
 R_0=\operatorname{im}
 \bigl(W\otimes U_3^\vee\longrightarrow U_0\otimes U_1\otimes U_2\bigr)
\]
を定める。geometric 条件によりこれは階数二の部分束である。
五つ組の四項 foundation 代数 $\Lambda(W)$ は、頂点 $0,1,2,3$、
隣接矢空間 $U_0,U_1,U_2$ を持ち、三次関係が $R_0$ である代数である。
従ってその前向き成分は
\[
 \Lambda(W)_{ii}=\OO_S,\quad
 \Lambda(W)_{i,i+1}=U_i,\quad
 \Lambda(W)_{i,i+2}=U_i\otimes U_{i+1},\quad
 \Lambda(W)_{03}=(U_0\otimes U_1\otimes U_2)/R_0.
\]
全階数は $24$ である。積は path の連結であり、圏の Hom と比較するときは合成による積を用いる。
同じ右加群規約を両辺で固定する。

普遍 cubic 代数 $D(w)$ を前向きの $\Hom(E_i,E_j)$ の順で読む場合、
右 tails の入力は、添字を反転した反対代数
\begin{equation}\label{eq:quintuple-right-tails-convention}
 C_{ij}=D(w)_{-j,-i},\qquad E_i=\pi e_{-i}C
\end{equation}
である。定理~\ref{thm:relative-cubic-tails-tilting} の全構造 Hom の計算により
\[
 \Hom(E_i,E_j)=C_{-j,-i}=D(w)_{ij}\qquad(i\le j)
\]
となり、上の foundation と一致する。
実際の二周期 $\beta_D:D\to D(2)$ に対しては
$\beta_C=\beta_D^{-1}:C\to C(2)$ を用いる。
この変換は反転した canonical 四周期との平方比較を保つ。
相対族への昇格では、定理~\ref{thm:geometric-cubic-relative-effectivity} をこの $C$ に適用する。
偶数環と右軌道環の比較は定理~\ref{thm:quadric-finite-type-effectivity} の証明と同じで、
右軌道積は $a*b=s^{\deg b}(a)\circ b$ である。
$D$ の通常の右 tails と $C$ の右 tails を同一視することはしない。

\subsection{五つの対象からの有限次数回復}

\begin{lemma}[四次 overlap の普遍的回復]
\label{lem:cubic-artin-overlap-recovery}
$B$ を剰余体 $k$ の可換局所 Artin $k$ 代数とし、
$\mathcal D$ を $B$ 線形 Grothendieck 圏とする。
$E_0,\ldots,E_4$ は compact な $B$-flat 対象であり、
その導来 Hom は前向きに次数零へ集中すると仮定する。
また $(E_0,E_1,E_2,E_3)$ と $(E_1,E_2,E_3,E_4)$ は strong exceptional collections とする。
\[
 D_{ij}=\Hom_{\mathcal D}(E_i,E_j)\quad(0\le i\le j\le4)
\]
は有限自由で、剰余体上への基底変更と可換し、
その剰余体上の積は標準 geometric cubic 代数の対応する窓と一致するとする。
このとき、次の構成が成立する。
\begin{enumerate}
\item $V_i=D_{i,i+1}$ は階数二である。二重積は同型、三重積と四重積は全射であり、
\[
 \operatorname{rk}D_{i,i+r}=2,4,6,9\qquad(r=1,2,3,4)
\]
である（添字が $0,\ldots,4$ に入る範囲で読む）。
\item 三重積の核
\[
 R_i=\ker\bigl(V_i\otimes V_{i+1}\otimes V_{i+2}\longrightarrow D_{i,i+3}\bigr),
 \qquad i=0,1,
\]
は階数二の直和因子である。
\item $T_4=V_0\otimes V_1\otimes V_2\otimes V_3$ と置く。
\[
 W_0=(R_0\otimes V_3)\cap(V_0\otimes R_1)\subset T_4
\]
は階数一の直和因子であり、積による列
\begin{equation}\label{eq:artin-overlap-complex}
0\longrightarrow W_0\longrightarrow
 (R_0\otimes V_3)\oplus(V_0\otimes R_1)
 \longrightarrow T_4\longrightarrow D_{04}\longrightarrow0
\end{equation}
は分裂完全な有限自由複体である。
最初の射は $w\mapsto(w,-w)$、中央の射は二つの包含の和である。
\item 縮約は同型
\begin{equation}\label{eq:artin-contraction-isos}
 W_0\otimes V_3^\vee\xrightarrow{\sim}R_0,
 \qquad
 W_0\otimes V_0^\vee\xrightarrow{\sim}R_1
\end{equation}
を与える。従って $(V_0,V_1,V_2,V_3,W_0)$ は $B$ 上の geometric 五つ組である。
\end{enumerate}
これらの有限加群、核、交わり、縮約の構成は、任意の可換 $B$ 代数への
基底変更と可換する。
圏の Hom の基底変更と比較するときは、与えられた相対 compact 圏の基底変更同型を用いる。
\end{lemma}

\begin{proof}
全証明は相対モジュライ論文 \cite[\texttt{cubic_smooth_atlas}]{QuadricModuli2026} に置く。
\end{proof}

\begin{corollary}[四項 foundation の完全な同定]
\label{cor:quintuple-four-window}
前補題で四つの基底同型 $U_i\otimes B\xrightarrow{\sim}V_i$ を選び、
それによる五つ組の $B$ 点を $w_B$ とする。
このとき合成は、頂点と全ての Hom の作用を保つ同型
\[
 \Lambda(w_B)\xrightarrow{\sim}
 \End_{\mathcal D}\left(\bigoplus_{i=0}^3E_i\right)
\]
を与える。
\end{corollary}

\begin{proof}
全証明は相対モジュライ論文 \cite[\texttt{cubic_smooth_atlas}]{QuadricModuli2026} に置く。
\end{proof}

\subsection{指定された還元同型と基底を持つ持上げ}

\begin{proposition}[一般 cubic 点での framed Artin 持上げ]
\label{prop:framed-cubic-artin-lift}
$A'\twoheadrightarrow A$ を剰余体 $k$ を持つ局所 Artin $k$ 代数の冪零拡大とする。
$\mathcal G_A$ は flat な $A$ 線形 Grothendieck 圏であり、
$A$-flat な対象の geometric $(4,3)$-helix $(E_{i,A})_{i\in\mathbb Z}$ を持つとする。
全ての前向き Hom は有限自由で基底変更と可換し、閉ファイバーは標準 geometric cubic helix であるとする。
特に各四項窓は $D(\mathcal G_A)$ の compact tilting generator である。

$\mathcal G_{A'}$ を $\mathcal G_A$ の flat Grothendieck 変形とする。
同時に指定された還元同型を持つ flat 構造対象 $\OO_{A'}$ があり、
$E_{0,A}=\OO_A$ と同定されているとする。
四つの基底同型
\[
 \varphi_{i,A}:U_i\otimes A\xrightarrow{\sim}
 \Hom(E_{i,A},E_{i+1,A}),\qquad i=0,1,2,3,
\]
を指定し、回復される五つ組の点を $w_A\in S(A)$ とする。
このとき次が存在する。
\begin{enumerate}
\item flat な geometric helix $(E'_{i})_{i\in\mathbb Z}$ と指定還元同型
  $A\otimes_{A'}E'_i\xrightarrow{\sim}E_{i,A}$。
  $E'_0=\OO_{A'}$ とでき、その還元同型も指定したものに取れる。
\item $\varphi_{i,A}$ を延長する基底同型
\[
 \varphi_{i,A'}:U_i\otimes A'\xrightarrow{\sim}\Hom(E'_i,E'_{i+1})
\]
と、これらから回復される $w_{A'}\in S(A')$。
その $A$ への還元は\emph{指定された点} $w_A$ に等しい。
\item 頂点と全ての積を保つ同型
\[
 \Lambda(w_{A'})\xrightarrow{\sim}\End\left(\bigoplus_{i=0}^3E'_i\right)
\]
であって、その還元は $\varphi_{i,A}$ と元の合成から定まる同型に一致するもの。
\end{enumerate}
$W_{0,A}$ の生成元も指定されていれば、その生成元を $W_{0,A'}$ の生成元へ持ち上げられる。
射影空間 $S$ の点の持上げはこの追加の生成元の選択を必要としない。
\end{proposition}

\begin{proof}
全証明は相対モジュライ論文 \cite[\texttt{cubic_smooth_atlas}]{QuadricModuli2026} に置く。
\end{proof}

\subsection{pointed exact 同型の還元を保つ比較}

\begin{lemma}[Artin 基底上の二つの heart の比較]
\label{lem:artin-two-heart-comparison}
$B$ は局所 Artin 環とし、$\mathcal G_1,\mathcal G_2$ を flat な $B$ 線形
Grothendieck 圏とする。それぞれ flat compact tilting generator $T_1,T_2$ を持ち、
その End 代数は $B$ 上有限自由で、Artin 商への基底変更と可換するとする。
各 Artin 商でもこれらの還元は compact tilting generators と仮定する。
頂点を保つ End 代数の同型から標準 dg Morita 同値
\[
 F:D(\mathcal G_1)\xrightarrow{\sim}D(\mathcal G_2)
\]
を定める。
その閉ファイバー $F_k$ が指定された exact 同値の導来化なら、$F$ も heart-exact である。
さらに $B\twoheadrightarrow A$ 上の還元 Morita 同値が
指定された exact 同値 $\alpha_A$ の導来化に一致すれば、
$F$ の heart への制限は $\alpha_A$ を延長する。
指定された構造対象とその比較同型を保つ場合は pointed 同値となる。
\end{lemma}

\begin{proof}
全証明は相対モジュライ論文 \cite[\texttt{cubic_smooth_atlas}]{QuadricModuli2026} に置く。
\end{proof}

\begin{theorem}[五つ組への étale chart の指定 Artin 持上げ]
\label{thm:quintuple-chart-artin-lifting}
$S\subset\mathscr G$ を上の開部分とし、$a:U\to S$ を有限型の étale 射とする。
実際の相対族 $q:U\to\mathfrak H_2$ が構成されており、任意の affine 基底変更に整合する
flat compact geometric cubic helix $(E_i)$、$E_0=\OO$ を持つとする。
さらに $a$ の普遍五つ組とこの helix の五つ組が同定され、
その四項窓の End 代数は頂点と積を保って $a^*\Lambda(W)$ と同定されるとする。
全ての幾何点で helix は標準 geometric cubic helix と仮定する。

$A'\twoheadrightarrow A$ を局所 Artin $k$ 代数の small extension とし、
$u_A\in U(A)$、$c_{A'}\in\mathfrak H_2(A')$、および\emph{指定された} pointed exact 同型
\[
 \alpha_A:q(u_A)\xrightarrow{\sim}c_{A'}|_A
\]
を取る。
このとき $u_A$ を延長する $u_{A'}\in U(A')$ と、$\alpha_A$ を延長する pointed exact 同型
\[
 \alpha_{A'}:q(u_{A'})\xrightarrow{\sim}c_{A'}
\]
が存在する。いずれの段階にも Hochschild 一次または二次上の根の作用の仮定は不要である。
\end{theorem}

\begin{proof}
全証明は相対モジュライ論文 \cite[\texttt{cubic_smooth_atlas}]{QuadricModuli2026} に置く。
\end{proof}

\subsection{smoothness の帰結と限界}

\begin{corollary}[全 closed cubic 点での smooth chart]
\label{cor:quintuple-effective-smooth-chart}
前定理の入力に加え、既存の相対認識から全忠実な射
\[
 \Phi:\mathfrak H_2\longrightarrow\mathfrak P=[X_3/\operatorname{PGL}_9]
\]
があり、$\mathfrak P$ の対角は affine かつ有限表示であるとする。
このとき $q:U\to\mathfrak H_2$ はスキームで表現可能な smooth 射である。
その fppf 像は $\mathfrak H_2$ の開部分スタックで、有限型で滑らかな Artin スタックである。
\end{corollary}

\begin{proof}
全証明は相対モジュライ論文 \cite[\texttt{cubic_smooth_atlas}]{QuadricModuli2026} に置く。
\end{proof}

\subsection{全 cubic locus の有限型 smooth 開部分}

\begin{theorem}[cubic 対象全体の smooth Artin 開部分]
\label{thm:cubic-smooth-open}
$\mathfrak H_2$ には有限型で滑らかな Artin 開部分
\[
 \mathfrak V_{\mathrm{cub}}\subset\mathfrak H_2
\]
が存在し、その幾何点はちょうど pointed cubic 表示を持つ対象である。
より正確には、任意の局所有限型基底 $T$ 上の相対 $\mathfrak H_2$ 対象が
$\mathfrak V_{\mathrm{cub}}(T)$ に属するための必要十分条件は、
全ての幾何ファイバーが pointed cubic 表示を持つことである。

有限型の滑らかなスキーム $U_{\mathrm{cub}}$ と representable smooth 全射
\[
 q_{\mathrm{cub}}:U_{\mathrm{cub}}\longrightarrow\mathfrak V_{\mathrm{cub}}
\]
を取れる。その同値関係は全ての pointed exact 同値を用いる
\[
 R_{\mathrm{cub}}=U_{\mathrm{cub}}\times_{\mathfrak H_2}U_{\mathrm{cub}}
 \simeq U_{\mathrm{cub}}\times_{\mathfrak P}U_{\mathrm{cub}},\qquad
 \mathfrak V_{\mathrm{cub}}\simeq[U_{\mathrm{cub}}/R_{\mathrm{cub}}]
\]
であり、選んだ五つ組や根を保つ同値だけに限定しない。
\end{theorem}

\begin{proof}
全証明は相対モジュライ論文 \cite[\texttt{cubic_smooth_atlas}]{QuadricModuli2026} に置く。
\end{proof}

\begin{corollary}[全 cubic 点の形式変形と次元]
\label{cor:cubic-formal-h2-abelian}
$x=(\mathcal C,\OO)$ を $\mathfrak V_{\mathrm{cub}}(k)$ の点とする。
剰余体 $k$ の局所 Artin 基底上で、閉ファイバーの同定を持つ相対 $\mathfrak H_2$ 変形から
flat pointed Abel 変形への忘却は groupoid の同値である。
従ってこの形式変形問題は無障害であり、四項 foundation 代数を $\Lambda$ とすると、
その接空間と無限小自己同型空間はそれぞれ
\[
 \operatorname{HH}^2(\Lambda),\qquad\operatorname{HH}^1(\Lambda)
\]
である。さらに $\mathfrak V_{\mathrm{cub}}$ は純次元三である。
ここでは根・根の平方同型・AS4 被覆は標識として固定しない。
\end{corollary}

\begin{proof}
全証明は相対モジュライ論文 \cite[\texttt{cubic_smooth_atlas}]{QuadricModuli2026} に置く。
\end{proof}

\begin{remark}[noncubic 境界との区別]
$\mathfrak V_{\mathrm{cub}}$ の幾何点は cubic 対象だけであり、
定理~\ref{thm:complete-h2-classification} の noncubic 対象は含まない。
従ってこの定理は $\mathfrak H_2$ 全体の代数性や、noncubic 境界に沿う相対認識を結論しない。
また $\mathfrak V_{\mathrm{cub}}$ を $[X_3/\operatorname{PGL}_9]$ の開部分と同一視する主張でもない。
結論に必要なのは、相対 Picard 作用と二周期の実際の構成、同じ根の相対 AS4 被覆、
全 Noetherian 性、canonical tails 比較、そして本節の指定 Artin 同型の持上げを合わせたことである。
\end{remark}
